\chapter{Konzeptentwurf für eine informationstechnische Anbindung der Komponenten} \label{cha:Konzeptentwurf}
Auf Basis der unter \ref{cha:Analyse} gewonnenen Erkenntnissen zum Speicherbedarf der Komponenten werden in diesem Kapitel Entwürfe erarbeitet, wie die Komponenten in eine Fahrzeugarchitektur implementiert werden können.
\section{Einteilung der Komponenten}
Im ersten Schritt werden die durch die Analyse gewonnenen Erkenntnisse genutzt, um Komponenten nach ihren Speichereigenschaften zu sortieren. \\
Die erste Gruppe benötigt einen Speicher für 50 Inhalte von weniger als $ 30\,\mathrm{MByte} $, die zweite Gruppe von weniger als $ 1\,\mathrm{GByte} $, die dritte Gruppe von weniger als $ 5\,\mathrm{GByte} $. \\
Tabelle \ref{tab:Einteilung} weist die Komponenten den Gruppen zu. Für die Konzeptentwürfe werden allen Komponenten die maximalen Speichergrößen der Gruppe zugeordnet. 
\\
\begin{table}[hbt]	
	\centering
	\renewcommand{\arraystretch}{1.5}	% Skaliert die Zeilenhöhe der Tabelle
	\captionabove[Einteilung der Komponenten nach Speichergröße]{Einteilung der Komponenten nach Speichergröße}
	\label{tab:Einteilung}
	\begin{tabular}{c|c|c}
		\textbf{Erste Gruppe (EG\nomenclature{EG}{Erste Gruppe})} & \textbf{Zweite Gruppe (ZG\nomenclature{ZG}{Zweite Gruppe})} & \textbf{Dritte Gruppe (DG\nomenclature{DG}{Dritte Gruppe})} \\ 
		$ 30\,\mathrm{MByte} $ & $ 1\,\mathrm{GByte} $ & $ 5\,\mathrm{GByte} $ \\
		\hline
		\hline
		\textbf{Bild-Komponenten} & \textbf{Bild-Komponenten} & \textbf{Bild-Komponenten} \\
		\parbox[t]{0.3\linewidth}{\centering E-Papier in der Frontschürze} &  &  \\
		\parbox[t]{0.3\linewidth}{\centering E-Papier über den vorderen Radkästen} &  & \\
		\parbox[t]{0.3\linewidth}{\centering E-Papier in der Heckleuchte} &  &  \\
		\textbf{Video-Komponenten} & \textbf{Video-Komponenten} & \textbf{Video-Komponenten} \\
		\parbox[t]{0.3\linewidth}{\centering LED-Streifen in der Frontschürze} & \parbox[t]{0.3\linewidth}{\centering LED-Matrix im Dachhimmel} & \parbox[t]{0.3\linewidth}{\centering Videoprojektoren im Fußraum} \\
		\parbox[t]{0.3\linewidth}{\centering LED-Streifen in den Radkästen} &  & \parbox[t]{0.3\linewidth}{\centering Videoprojektoren in\\den Außenspiegeln} \\ 
		\parbox[t]{0.3\linewidth}{\centering LED-Streifen in der Heckleuchte} &  & \parbox[t]{0.3\linewidth}{\centering Bildschirme in den\\hinteren Seitenfenstern} \\ 
		\parbox[t]{0.3\linewidth}{\centering LED-Streifen im Interieur} &  & \parbox[t]{0.3\linewidth}{\centering Bildschirme in den\\hinteren Seitenfenstern} \\
		\parbox[t]{0.3\linewidth}{\centering LED Türtafeln} & & \parbox[t]{0.3\linewidth}{\centering  Bildschirme in der Einstiegsleiste} \\
		\parbox[t]{0.3\linewidth}{\centering Morphende Oberfläche\\in der Mittelkonsole} & & \parbox[t]{0.3\linewidth}{\centering Durchsichtiger Bildschirm\\im Dachfenster} \\
	\end{tabular} 
\end{table}
Diese erste Einteilung liefert eine Gruppierung mit der Speichergröße. Eine weitere Einteilung nach Komponenten, die statisch Bilder anzeigen oder dynamisch Videos anzeigen, ist relevant für die Dimensionierung der Busanbindung. \\
Bei Komponenten beträgt die Größe eines Inhalts die Speichergröße der Gruppe durch die 50 Inhalte. Bei Komponenten, die nur Bilder anzeigen ist dieser Betrag gleich des Betrags für ein Bild.
\begin{align}
	EG &= \frac{30\,\mathrm{MByte}}{50} = 600\,\mathrm{kByte} \\
	ZG &= \frac{1\,\mathrm{GByte}}{50} =  20\,\mathrm{MByte} \\
	DG &= \frac{5\,\mathrm{GByte}}{50} = 100\,\mathrm{MByte}
\end{align}
Bei Komponenten, die Videos anzeigen, beträgt die Größe der einzelnen Bilder in den Videos die Größe der oben berechneten Inhalte durch die Anzahl der Bilder pro Videos.
\begin{align}
	EG &= \frac{600\,\mathrm{kByte}}{10\,\mathrm{s} \cdot 24\,\mathrm{fps}} = 2,50\,\mathrm{kByte} \\
	ZG &= \frac{20\,\mathrm{MByte}}{10\,\mathrm{s} \cdot 24\,\mathrm{fps}} = 83,33\,\mathrm{kByte} \\
	DG &= \frac{100\,\mathrm{MByte}}{10\,\mathrm{s} \cdot 24\,\mathrm{fps}} = 416,67\,\mathrm{kByte}
\end{align}
Die Datenrate der Video-Komponenten können berechnet werden, indem die Größe der Inhalt durch die Länge des Videos dividiert wird.
\begin{align}
	EG &= \frac{600\,\mathrm{kByte}}{10\,\mathrm{s}} = 60\,\frac{\mathrm{kByte}}{\mathrm{s}}\\
	ZG &= \frac{20\,\mathrm{MByte}}{10\,\mathrm{s}} = 2.000\,\frac{\mathrm{kByte}}{\mathrm{s}} \\
	DG &= \frac{100\,\mathrm{MByte}}{10\,\mathrm{s}} = 10.000\,\frac{\mathrm{kByte}}{\mathrm{s}}
\end{align}
Die erste Gruppe besteht aus Bild- und Videokomponenten, die zweite und dritte Gruppe nur aus Videokomponenten.


%TODO Ab Hier weitermachen%

Abweichend von der Betrachtungsweise der verwendeten Technologie und des Einbauortes, ist es sinnvoll nach den Art der Informationsbedarfes zu unterteilen. Die Unterteilung 
Die Anforderungen an die Komponenten sind gleich. Jede Komponente muss in Echtzeit auf mindestens 10 Kollektionen umschalten können.
\section{Lösungsansätze unterschiedlicher Konzeptentwurf}
Eine Komponente besteht in diesem Fall immer aus einer Anzeigefläche mit einer dazugehörigen Steuerung, die aus Bild oder Videodateien die einzelnen Bildpunkte setzt.

Die Kollektionen werden über eine Drahtlose Verbindung zwischen einem zentralen Steuergerät aus dem Internet geladen und müssen auf den Komponenten im Fahrzeug dargestellt werden. 
Das zentrale Abspeichern der Daten und anzeigen der Daten bei Auswahl ist eine mögliche Option.
Die Daten können auch über das Fahrzeugnetzwerk vom zentralen Steuergerät auf verteilte Steuergeräte, die direkt an den Komponenten sind, geladen werden, wenn die Daten zu dem Zeitpunkt nicht angezeigt werden sollen. Bei Aktivierung der Kollektion müssen dann nur Steuersignale den Steuergeräten die Information zur Aussendung der Kollektion geben.
Relevant für die möglichen Optionen, sind die Latenz eines Steuersignals und der Dauer eines kompletten Downloads der Daten.
Der Anschluss der Komponenten an eine bestimmte Architektur ist unter unterschiedlichen Gesichtspunkten zu bewerten.

Es gibt unterschiedliche Varianten die Geräte anzuschliessen.

Die erste Form von Komponenten sind diejenigen, die nicht dafür ausgelegt sind neue Inhalte in Echtzeit vom zentralen System herunterzuladen.
Die zweite Form sind Komponenten, die in Echtzeit neu Inhalte vom System herunterladen müssen und direkt anzeigen.

Echtzeit bedeutet in diesem Fall, dass das System neue Inhalte anzeigt, die der Benutzer über das Zentraldisplay oder Drittgerät ausgewählt hat.
Durch diese Unterscheidung kann die zur Verfügung gestellte Bandbreite deutlich unterschiedlich sein. 
Ein weiterer wichtiges Merkmal ist die Geschwindigkeit zum 


Architekturdiagramm einer Möglichen Variante:
\begin{figure}[hbt]
	\centering
	\begin{tikzpicture}
	\node[text width=3.2cm, align=center] (Computer) {Computer};
\end{tikzpicture}
	\caption[Netzwerkarchitektur Ansteuerung Komponenten]{Netzwerkarchitektur Ansteuerung Komponenten}
	\label{fig:architektur1}
\end{figure}
\section{Vergleich der unterschiedlichen Lösungsansätze}





